% =====================================================================
%  Chapter 31 — The Worked Cases
% =====================================================================

\chapter{The Worked Cases}
\label{ch:worked-cases}

A rule is worth nothing until it has been run. This chapter runs the
selection rule of Chapter 30 on the two documented cases, step by step, and
shows the numbers that come out at each step.

\section{Case one: 2024 YR4 at the 2032 lunar encounter}

The inputs are the published ones: an effective diameter of 60.8 metres with
an uncertainty of \(+1.8\) and \(-3.6\), an absolute magnitude of 23.93 from
the Small-Body Database, a heliocentric distance near 1.1 astronomical units,
a bulk density bracketed between 1,300 and 3,500 kilogrammes per cubic metre,
and a close approach on 22 December 2032 at a lunar distance of 22,900
kilometres.

\textbf{Step 1, the required baseline.} The angular diameter at 22,900
kilometres is 548 milliarcseconds. For four samples across the diameter at
750 nanometres, the matched baseline is 1.13 metres; for seven samples it is
1.98 metres. The aperture that resolves the disk outright is 0.345 metres.

\textbf{Step 2, the over-resolution limit.} The half-visibility baseline is
0.40 metres, which is below the matched baseline. This is not a wall, and
the reason is worth stating because it is easy to get wrong: the
half-visibility criterion is the limit of the \emph{centroid} mode, in which
an unresolved source is located from the centre of the fringe pattern. A
resolved body is imaged from the visibility function itself, out to the
baseline where the signal is lost in the noise, and there the practical limit
is the signal-to-noise rather than a visibility threshold. It is worth
noticing that the uniform-disc visibility at the matched baseline is 0.068
in this case and in every other, because the matched baseline is by
definition the one at which the object's diameter spans four fringe
spacings, and the visibility at that geometry is a constant.

\textbf{Step 3, the astrometry.} The centroid precision at the matched
baseline, with a signal-to-noise ratio of 7, is 19.6 milliarcseconds, which
at 22,900 kilometres is a position of 2.17 metres. The lunar impact
probability is set by the object's position on the Moon's surface, and the
Moon's surface is thousands of kilometres across. A 2.17-metre position is
vastly sufficient.

And the range is not inferred at all: the observer's position is known to
centimetres, so the range to an object at 22,900 kilometres is known to
centimetres.

\textbf{Step 4, buildability.} A 1.13-metre baseline is buildable on the
Earth, in a laboratory. The interesting versions of this case are the ones
where the object is not at a lunar distance.

\textbf{Step 5, the window.} The object's brightness window at that
distance is weeks.

\textbf{Step 6, the condition.} This is where the case is decided, and it is
decided against the lunar platform. The astrometric B-plane uncertainty for
the 2032 encounter was 270,000 kilometres, reducible to 25,000 kilometres or
less by a single extra observation, which the ground would achieve by the
end of the apparition. Against that, the along-track displacement produced
by a thirty per cent error in the Yarkovsky drift rate is 17.2 kilometres
over seven years in the fastest case tabulated.

The astrometric term dominates by three orders of magnitude. The rule
declines the lunar platform for this target, and it should.

\begin{table}[!htbp]
\centering
\small
\caption{Case one: the rule applied to 2024 YR4 at the 2032 lunar
encounter.}
\label{tab:case1}
\begin{tabularx}{\textwidth}{@{}P{0.61in}P{1.95in}Y@{}}
\toprule
\rowcolor{panel}
\textbf{Step} & \textbf{Result} & \textbf{Reading} \\
\midrule
1 & baseline 1.13 m; aperture 0.345 m & the object is resolved at a lunar
distance \\
\rowcolor{panel}
2 & half-visibility baseline 0.40 m & beyond it, in the imaging mode; not a
wall \\
3 & 19.6 mas, or 2.17 m at the range & and the range is known directly \\
\rowcolor{panel}
4 & buildable & a short baseline is not the lunar argument \\
5 & a window of weeks & no scheduling problem \\
6 & astrometric term 3 orders larger & \textbf{decline the lunar platform} \\
\bottomrule
\end{tabularx}
\end{table}

\section{Case two: Apophis at the 2029 Earth encounter}

The inputs: a diameter of 325 metres, an absolute magnitude of 19.09, and a
close approach on 13 April 2029 at 38,011 kilometres with a relative speed
of 7.42 kilometres per second.

\textbf{Step 1.} The angular diameter is 1.76 arcseconds. The matched
baseline for four samples is 0.35 metres, and the aperture that resolves
the disk is 0.107 metres.

\textbf{Step 2.} The half-visibility baseline is 0.124 metres, again below
the matched baseline, for the same reason.

\textbf{Step 3.} The centroid precision at the matched baseline is 63.0
milliarcseconds, which at 38,011 kilometres is a position of 11.6 metres.
The range is direct and known to the centimetre, because the observer is on
the Moon and the Moon's position is known.

\textbf{Step 4.} A 0.35-metre baseline is buildable anywhere. The lunar
argument does not rest on this case either.

\textbf{Step 5.} The brightness window is a few days around the encounter.

\textbf{Step 6.} And here the rule selects. The along-track standard
deviation for the 2036 encounter is 398.5 kilometres. The diameter
assumption alone is worth 507 kilometres of keyhole displacement, and the
density assumption 379 kilometres. Each of them is comparable with the
astrometric term, and neither of them is reduced by any quantity of
additional astrometry.

\begin{table}[!htbp]
\centering
\small
\caption{Case two: the rule applied to Apophis at the 2029 encounter.}
\label{tab:case2}
\begin{tabularx}{\textwidth}{@{}P{0.61in}P{1.95in}Y@{}}
\toprule
\rowcolor{panel}
\textbf{Step} & \textbf{Result} & \textbf{Reading} \\
\midrule
1 & baseline 0.35 m; aperture 0.107 m & the object is resolved at a
ten-thousand-kilometre range \\
\rowcolor{panel}
2 & half-visibility baseline 0.124 m & imaging mode again \\
3 & 63.0 mas, or 11.6 m at the range & the range is direct \\
\rowcolor{panel}
4 & buildable & not the lunar argument \\
5 & a window of days & schedulable \\
6 & physical terms comparable with the astrometric sigma & \textbf{select
the lunar platform} \\
\bottomrule
\end{tabularx}
\end{table}

\section{What the two cases together establish}


Read together, the two cases say something that neither says alone.

The measurement capability at a close encounter is not the scarce thing. A
matched baseline of a metre, an aperture of a third of a metre and a
position of a few metres are all obtainable. The scarce thing is the
\emph{presence} at the encounter, and the presence is what the lunar
platform offers.

And the presence is worth what the physical terms are worth --- 507
kilometres of keyhole displacement in the 2036 case, and 379 for the
density --- and not one metre more.

\begin{seamnote}{The book's own prediction, stated so it can be checked}
For 2024 YR4, the book predicts that the resolved shape measured during the
2032 lunar encounter will fix the volume, and that the energy implied will
fall inside the published model set's range of 2.16 to 9.79 megatons rather
than outside it. If the measured shape instead implies an energy outside
that range, the book's account of the thermal-model degeneracy is wrong and
the reason will be in the shape. That is the prediction that would fail
first, and it is stated here so that it can be.
\end{seamnote}

\begin{keytakeaways}
\item Both cases are resolved at the encounter range with a matched baseline
of a metre or less and an aperture of a third of a metre, and the astrometric
position is a few metres.
\item At the 2024 YR4 encounter the astrometric term dominates by three
orders of magnitude and the rule declines the lunar platform.
\item At the Apophis encounter the physical terms are comparable with the
astrometric standard deviation and the rule selects it.
\item The scarce thing is not the measurement but the presence at the
encounter, and the presence is worth the physical term and no more.
\end{keytakeaways}
